\chapter{Teoría y Fundamentos de Probabilidad}
\label{cap:probabilidades}

\section*{Objetivos de Aprendizaje de la Unidad}
Al finalizar el estudio de esta unidad, el estudiante será capaz de:
\begin{enumerate}
    \item Comprender los fundamentos de la incertidumbre en los negocios y diferenciar entre los enfoques \textbf{clásico}, de \textbf{frecuencia relativa} y \textbf{subjetivo} de probabilidad.
    \item Aplicar con exactitud los axiomas de probabilidad y las reglas de la \textbf{adición} y de la \textbf{multiplicación} para eventos dependientes e independientes.
    \item Construir e interpretar \textbf{tablas de contingencia bidimensionales} para calcular probabilidades conjuntas, marginales y condicionales en análisis de mercado y auditoría.
    \item Dominar y aplicar el \textbf{Teorema de Bayes} para actualizar probabilidades a priori con nueva información en análisis de riesgo crediticio, fraude contable y toma de decisiones gerenciales.
\end{enumerate}

\section{Introducción a la Probabilidad en el Entorno Contable-Empresarial}
En el mundo de los negocios, la contabilidad y las finanzas, casi ninguna decisión directiva se toma bajo condiciones de certeza absoluta. Un auditor no revisa el 100\% de las transacciones de una empresa, sino una muestra; un oficial de crédito no puede predecir con certeza si un prestatario pagará a tiempo; y un analista financiero debe evaluar proyectos con flujos de caja sujetos a riesgo de mercado.

La \textbf{Probabilidad} es el lenguaje matemático que permite medir y modelar cuantitativamente la incertidumbre, permitiendo transformar corazonadas subjetivas en decisiones fundamentadas.

\section{Conceptos Fundamentales}

\begin{definicion}{Experimento, Espacio Muestral y Eventos}
\begin{itemize}
    \item \textbf{Experimento Aleatorio:} Proceso mediante el cual se genera una observación o medición, cuyo resultado exacto no se puede predecir con certeza de antemano (ej. auditar una factura para ver si tiene o no error tributario).
    \item \textbf{Espacio Muestral ($S$):} Conjunto de todos los resultados elementales posibles de un experimento aleatorio.
    \item \textbf{Punto Muestral ($e_i$):} Cada uno de los resultados elementales e indivisibles que conforman el espacio muestral.
    \item \textbf{Evento o Suceso ($A$):} Cualquier subconjunto del espacio muestral ($A \subseteq S$).
\end{itemize}
\end{definicion}

\section{Enfoques Conceptuales de la Probabilidad}

\subsection{1. Enfoque Clásico o \textit{A Priori} (Laplace)}
Se aplica cuando todos los resultados elementales del espacio muestral son \textbf{equiprobables} (tienen la misma posibilidad de ocurrir).
\begin{equation}
    P(A) = \frac{N(A)}{N(S)} = \frac{\text{Número de resultados favorables al evento } A}{\text{Número total de resultados posibles en } S}
\end{equation}

\subsection{2. Enfoque de Frecuencia Relativa o Empírico}
Se basa en la observación histórica de la proporción de veces que ocurre un evento al repetir el experimento un número grande $n$ de veces bajo condiciones similares:
\begin{equation}
    P(A) \approx \frac{n(A)}{n} = \frac{\text{Número de veces que ocurrió } A}{\text{Número total de ensayos realizados}}
\end{equation}
Por la \textbf{Ley de los Grandes Números}, a medida que el número de ensayos $n$ tiende a infinito, la frecuencia relativa converge a la verdadera probabilidad teórica.

\subsection{3. Enfoque Subjetivo o Personalista}
Se utiliza cuando no se dispone de datos históricos repetitivos ni de resultados equiprobables. Representa el \textbf{grado de creencia o confianza} que un decisor u experto asigna a la ocurrencia de un evento, con base en su experiencia, información cualitativa y juicio profesional.

\section{Axiomas de Kolmogorov y Propiedades Fundamentales}
Para cualquier evento $A$ definido en un espacio muestral $S$:
\begin{enumerate}
    \item \textbf{Axioma de No Negatividad:} $0 \le P(A) \le 1$ para todo evento $A$.
    \item \textbf{Axioma del Suceso Seguro:} $P(S) = 1$.
    \item \textbf{Suceso Imposible:} $P(\emptyset) = 0$.
    \item \textbf{Regla del Complemento:} La suma de la probabilidad de que ocurra $A$ más la de que no ocurra ($A'$) es igual a 1:
    \begin{equation}
        P(A) + P(A') = 1 \implies P(A') = 1 - P(A)
    \end{equation}
\end{enumerate}

\section{Reglas de la Adición de Probabilidades}

\subsection{Regla Especial de la Adición (Eventos Mutuamente Excluyentes)}
Dos eventos $A$ y $B$ son \textbf{mutuamente excluyentes} (o disjuntos) si no pueden ocurrir simultáneamente ($A \cap B = \emptyset$).
\begin{equation}
    P(A \cup B) = P(A) + P(B)
\end{equation}

\subsection{Regla General de la Adición (Eventos No Excluyentes)}
Si los eventos $A$ y $B$ pueden ocurrir juntos ($A \cap B \neq \emptyset$):
\begin{equation}
    P(A \cup B) = P(A) + P(B) - P(A \cap B)
\end{equation}
donde $P(A \cap B)$ representa la probabilidad conjunta de que ocurran ambos eventos.

\begin{ejemplo}{Auditoría de Cumplimiento Tributario (Basado en Anderson)}
Una firma auditora en Santa Cruz revisa los registros fiscales de 200 medianas empresas. Se encontró que:
\begin{itemize}
    \item El $40\%$ de las empresas presenta inconsistencias en el Libro de Compras IVA (Evento $A$).
    \item El $25\%$ presenta inconsistencias en el Impuesto a las Utilidades IUE (Evento $B$).
    \item El $15\%$ presenta inconsistencias en \textbf{ambos} tributos ($A \cap B$).
\end{itemize}
Si se selecciona una empresa al azar, ¿cuál es la probabilidad de que tenga problemas tributarios en el IVA o en el IUE (o en ambos)?

\tcbline
\textbf{Solución:}
Identificamos las probabilidades dadas:
\[
P(A) = 0.40, \quad P(B) = 0.25, \quad P(A \cap B) = 0.15
\]
Aplicando la Regla General de la Adición:
\[
P(A \cup B) = P(A) + P(B) - P(A \cap B) = 0.40 + 0.25 - 0.15 = 0.50 \implies 50\%
\]
\textbf{Interpretación:} La probabilidad de que una empresa presente al menos una observación fiscal en IVA o IUE es del 50\%.
\end{ejemplo}

\section{Probabilidad Condicional e Independencia}

\subsection{Definición de Probabilidad Condicional}
\begin{definicion}{Probabilidad Condicional}
La probabilidad condicional del evento $A$, dado que se sabe con certeza que ya ocurrió el evento $B$ (con $P(B) > 0$), se define formalmente como:
\begin{equation}
    P(A|B) = \frac{P(A \cap B)}{P(B)}
\end{equation}
De manera análoga, la probabilidad de $B$ dado $A$ es:
\begin{equation}
    P(B|A) = \frac{P(A \cap B)}{P(A)}, \quad \text{con } P(A) > 0
\end{equation}
\end{definicion}

\subsection{Independencia Estadística de Eventos}
\begin{definicion}{Eventos Independientes}
Dos eventos $A$ y $B$ son estadísticamente independientes si la ocurrencia de uno no afecta en absoluto la probabilidad de ocurrencia del otro:
\begin{equation}
    P(A|B) = P(A) \quad \text{o equivalentemente} \quad P(B|A) = P(B)
\end{equation}
Si esta igualdad no se cumple, se dice que los eventos son \textbf{dependientes}.
\end{definicion}

\section{Reglas de la Multiplicación}

\subsection{Regla Especial de la Multiplicación (Eventos Independientes)}
Si $A$ y $B$ son independientes:
\begin{equation}
    P(A \cap B) = P(A) \cdot P(B)
\end{equation}

\subsection{Regla General de la Multiplicación (Eventos Dependientes)}
Para cualquier par de eventos:
\begin{equation}
    P(A \cap B) = P(A) \cdot P(B|A) = P(B) \cdot P(A|B)
\end{equation}

\section{Tablas de Contingencia y Probabilidad Conjunta}
Una tabla de contingencia clasifica una muestra o población respecto a dos variables categóricas simultáneas.

\begin{ejemplo}{Riesgo Crediticio y Tamaño de Empresa (Basado en Lind)}
El departamento de riesgo crediticio de un banco comercial analizó el historial de 500 créditos corporativos otorgados durante el último año, arrojando los siguientes datos:

\begin{center}
\begin{tabular}{lccc}
\toprule
\textbf{Categoría de Crédito} & \textbf{Microempresa ($M$)} & \textbf{Gran Empresa ($G$)} & \textbf{Total} \\
\midrule
Al Día / Cumplido ($C$) & 260 & 190 & \textbf{450} \\
En Mora / Incumplido ($D$) & 40 & 10 & \textbf{50} \\
\midrule
\textbf{Total} & \textbf{300} & \textbf{200} & \textbf{500} \\
\bottomrule
\end{tabular}
\end{center}

Si se elige un crédito al azar, determine:
\begin{enumerate}[label=\alph*)]
    \item La probabilidad de que el crédito corresponda a una Microempresa: $P(M)$.
    \item La probabilidad conjunta de que sea de una Microempresa y esté en Mora: $P(M \cap D)$.
    \item La probabilidad de que esté en Mora, dado que es de una Microempresa: $P(D|M)$.
    \item La probabilidad de que esté en Mora, dado que es de una Gran Empresa: $P(D|G)$.
    \item ¿Es el estado del crédito independiente del tamaño de la empresa?
\end{enumerate}

\tcbline
\textbf{Solución:}
\begin{enumerate}[label=\alph*)]
    \item $P(M) = \frac{300}{500} = 0.60 \implies 60\%$.
    \item $P(M \cap D) = \frac{40}{500} = 0.08 \implies 8\%$.
    \item $P(D|M) = \frac{P(M \cap D)}{P(M)} = \frac{40/500}{300/500} = \frac{40}{300} = 0.1333 \implies 13.33\%$.
    \item $P(D|G) = \frac{P(G \cap D)}{P(G)} = \frac{10/500}{200/500} = \frac{10}{200} = 0.0500 \implies 5.00\%$.
    \item \textbf{Evaluación de Independencia:}
    La probabilidad global de mora es $P(D) = \frac{50}{500} = 0.10$ ($10\%$).\\
    Dado que $P(D|M) = 13.33\% \neq P(D) = 10\%$, el tamaño de la empresa y la condición de mora \textbf{son eventos dependientes}. Las microempresas presentan una tasa de mora significativamente superior a la de las grandes corporaciones.
\end{enumerate}
\end{ejemplo}

\section{Teorema de la Probabilidad Total y Teorema de Bayes}

\subsection{Teorema de la Probabilidad Total}
Si los eventos $A_1, A_2, \dots, A_k$ forman una \textbf{partición} del espacio muestral $S$ (son mutuamente excluyentes dos a dos y su unión conforma todo $S$, con $P(A_i) > 0$), entonces para cualquier evento arbitrario $B \subseteq S$:
\begin{equation}
    P(B) = \sum_{i=1}^k P(A_i) \cdot P(B|A_i) = P(A_1)P(B|A_1) + P(A_2)P(B|A_2) + \dots + P(A_k)P(B|A_k)
\end{equation}

\subsection{Teorema de Bayes}
\begin{definicion}{Teorema de Bayes}
Bajo las mismas condiciones de partición, la probabilidad condicional \textit{a posteriori} de un evento causante $A_j$, dado que ya se ha observado la evidencia $B$, está dada por:
\begin{equation}
    P(A_j|B) = \frac{P(A_j) \cdot P(B|A_j)}{\sum_{i=1}^k P(A_i) \cdot P(B|A_i)}
\end{equation}
\end{definicion}

\begin{ejemplo}{Detección de Fraude en Transacciones Electrónicas (Basado en Anderson)}
Un sistema contable centralizado procesa órdenes de pago de tres departamentos: Ventas ($A_1$), Adquisiciones ($A_2$) y Logística ($A_3$).
\begin{itemize}
    \item Ventas genera el $50\%$ de las órdenes ($P(A_1) = 0.50$).
    \item Adquisiciones genera el $30\%$ ($P(A_2) = 0.30$).
    \item Logística genera el $20\%$ ($P(A_3) = 0.20$).
\end{itemize}
La auditoría interna ha detectado que la probabilidad de que una transacción presente irregularidades contables o sospecha de fraude (Evento $F$) según su departamento de origen es:
\[
P(F|A_1) = 0.01 \quad (1\%), \qquad P(F|A_2) = 0.04 \quad (4\%), \qquad P(F|A_3) = 0.02 \quad (2\%)
\]
Durante un control de rutina, el sistema de alertas tempranas detecta una transacción con irregularidades graves. ¿Cuál es la probabilidad de que dicha transacción fraudulenta provenga del departamento de Adquisiciones ($A_2$)?

\tcbline
\textbf{Solución Paso a Paso:}
\begin{enumerate}
    \item \textbf{Cálculo de la probabilidad global de irregularidad $P(F)$ (Probabilidad Total):}
    \begin{align*}
        P(F) &= P(A_1)P(F|A_1) + P(A_2)P(F|A_2) + P(A_3)P(F|A_3) \\
        &= (0.50)(0.01) + (0.30)(0.04) + (0.20)(0.02) \\
        &= 0.005 + 0.012 + 0.004 = 0.021 \quad (2.1\%)
    \end{align*}
    \item \textbf{Aplicación del Teorema de Bayes para calcular $P(A_2|F)$:}
    \[
    P(A_2|F) = \frac{P(A_2)P(F|A_2)}{P(F)} = \frac{(0.30)(0.04)}{0.021} = \frac{0.012}{0.021} \approx 0.5714 \implies 57.14\%
    \]
\end{enumerate}
\textbf{Interpretación Contable:} Aunque el departamento de Adquisiciones solo origina el 30\% del volumen total de transacciones (probabilidad a priori), si se detecta una irregularidad, la probabilidad de que provenga de Adquisiciones se eleva drásticamente al \textbf{57.14\%} (probabilidad a posteriori). La auditoría forense debe priorizar la revisión en dicha área.
\end{ejemplo}

\section{Formulario Resumen del Capítulo 2}
\begin{formulario}[Resumen de Fórmulas — Teoría de Probabilidades]
\begin{itemize}
    \item \textbf{Probabilidad Clásica:} $P(A) = \dfrac{N(A)}{N(S)}$
    \item \textbf{Regla del Complemento:} $P(A') = 1 - P(A)$
    \item \textbf{Regla General de la Adición:} $P(A \cup B) = P(A) + P(B) - P(A \cap B)$
    \item \textbf{Probabilidad Condicional:} $P(A|B) = \dfrac{P(A \cap B)}{P(B)}$
    \item \textbf{Regla de Multiplicación (General):} $P(A \cap B) = P(A) \cdot P(B|A)$
    \item \textbf{Regla de Multiplicación (Independientes):} $P(A \cap B) = P(A) \cdot P(B)$
    \item \textbf{Teorema de Bayes:} $P(A_j|B) = \dfrac{P(A_j) \cdot P(B|A_j)}{\sum_{i=1}^k P(A_i) \cdot P(B|A_i)}$
\end{itemize}
\end{formulario}

\section{Ejercicios Propuestos de la Unidad}

\subsection*{Nivel 1: Conceptos y Reglas de Probabilidad}
\begin{enumerate}
    \item Enuncie los tres axiomas fundamentales de probabilidad de Kolmogorov y explique su significado en términos de certeza e imposibilidad.
    \item Sean $A$ y $B$ dos eventos asociados a un experimento aleatorio tales que $P(A) = 0.50$, $P(B) = 0.30$ y $P(A \cap B) = 0.15$.
    \begin{enumerate}[label=\alph*)]
        \item Calcule la probabilidad de la unión $P(A \cup B)$.
        \item Calcule la probabilidad condicional $P(A|B)$.
        \item Determine si los eventos $A$ y $B$ son estadísticamente independientes. Justifique formalmente.
    \end{enumerate}
    \item Si $P(E_1) = 0.40$, $P(E_2) = 0.35$ y $E_1$ y $E_2$ son mutuamente excluyentes, calcule:
    \begin{enumerate}[label=\alph*)]
        \item $P(E_1 \cup E_2)$
        \item $P(E_1 \cap E_2)$
        \item $P(E_1|E_2)$
    \end{enumerate}
\end{enumerate}

\subsection*{Nivel 2: Tablas de Contingencia y Probabilidad Condicional}
\begin{enumerate}[resume]
    \item Una firma de auditoría clasificó una muestra de 200 empresas según su tamaño (Pequeña, Mediana, Grande) y su cumplimiento en la presentación de estados financieros a tiempo (A tiempo, Con retraso):
    \begin{center}
    \small
    \begin{tabular}{l|cc|c}
    \toprule
    \textbf{Tamaño de Empresa} & \textbf{A Tiempo ($T$)} & \textbf{Con Retraso ($R$)} & \textbf{Total} \\
    \midrule
    Pequeña ($P$) & 60 & 40 & 100 \\
    Mediana ($M$) & 45 & 15 & 60 \\
    Grande ($G$) & 35 & 5 & 40 \\
    \midrule
    \textbf{Total} & \textbf{140} & \textbf{60} & \textbf{200} \\
    \bottomrule
    \end{tabular}
    \end{center}
    Si se selecciona una empresa al azar:
    \begin{enumerate}[label=\alph*)]
        \item Calcule la probabilidad de que haya presentado sus estados financieros con retraso.
        \item Calcule la probabilidad condicional de que presente con retraso dado que es una empresa pequeña: $P(R|P)$.
        \item Calcule la probabilidad condicional de que sea una empresa grande dado que presentó a tiempo: $P(G|T)$.
        \item ¿Es el cumplimiento tributario independiente del tamaño de la empresa? Justifique.
    \end{enumerate}
\end{enumerate}

\subsection*{Nivel 3: Teorema de Bayes y Toma de Decisiones}
\begin{enumerate}[resume]
    \item En una compañía comercializadora de granos en Montero, el 70\% de los clientes paga con transferencia bancaria electrónica y el 30\% con cheque. De los pagos efectuados con cheque, el 10\% son rechazados por falta de fondos o firmas disconformes. De los pagos por transferencia, solo el 1\% presenta inconsistencias operativas. Si el departamento de cobranzas detecta un pago con problemas:
    \begin{enumerate}[label=\alph*)]
        \item ¿Cuál es la probabilidad global de que un pago elegido al azar tenga problemas?
        \item ¿Cuál es la probabilidad de que el pago observado haya sido efectuado mediante cheque?
    \end{enumerate}
    \item Una entidad bancaria clasifica a sus solicitantes de crédito hipotecario en tres categorías de riesgo: Bajo ($60\%$), Medio ($30\%$) y Alto ($10\%$). La probabilidad histórica de entrar en mora prolongada es del 1\% para la categoría de riesgo bajo, 5\% para riesgo medio y 20\% para riesgo alto. Si un cliente incurre en mora prolongada, determine la probabilidad a posteriori de que pertenezca a la categoría de riesgo alto.
\end{enumerate}
